% Unified_Structural_Foundation.tex
% Title: Unified Structural Foundations of SuperBEST: Four Lower-Bound Techniques
% Author: Arturo R. Almaguer
% Date: April 2026
% Status: THEOREM NOTE (self-contained short note, 3-4 pages)

\documentclass[12pt,a4paper]{article}
\usepackage{amsmath,amssymb,amsthm}
\usepackage{geometry}
\usepackage{booktabs}
\usepackage{hyperref}
\geometry{margin=2.5cm}

%% ── Theorem environments ──────────────────────────────────────────────────────
\newtheorem{theorem}{Theorem}[section]
\newtheorem{lemma}[theorem]{Lemma}
\newtheorem{corollary}[theorem]{Corollary}
\newtheorem{proposition}[theorem]{Proposition}
\theoremstyle{definition}
\newtheorem{definition}[theorem]{Definition}
\newtheorem{remark}[theorem]{Remark}

%% ── Operator macros ───────────────────────────────────────────────────────────
\newcommand{\EML}{\operatorname{EML}}
\newcommand{\EXL}{\operatorname{EXL}}
\newcommand{\ELAd}{\operatorname{ELAd}}
\newcommand{\ELSb}{\operatorname{ELSb}}
\newcommand{\LEdiv}{\operatorname{LEdiv}}

%% ── Symbol macros ─────────────────────────────────────────────────────────────
\newcommand{\F}{\mathcal{F}_{16}}
\newcommand{\SB}{\mathrm{SB}}
\newcommand{\recip}{\operatorname{recip}}
\newcommand{\addgen}{\operatorname{add}_{\mathrm{gen}}}

%% =============================================================================
\title{Unified Structural Foundations of SuperBEST:\\
       Four Lower-Bound Techniques}
\author{Arturo R.\ Almaguer\\
  \small monogate research\quad\texttt{almaguer1986@gmail.com}}
\date{April 2026}
%% =============================================================================

\begin{document}
\maketitle

%% ── Abstract ──────────────────────────────────────────────────────────────────
\begin{abstract}
The SuperBEST routing table records minimum $\mathcal{F}_{16}$-node counts
for standard arithmetic operations.  Proving each entry optimal requires a
lower-bound argument.  We identify and formalize four complementary structural
techniques: T35 (non-mergeability of operator branches), T36 (tree-counting
argument), T37 (disjointness of reachable value sets), and T40 (linear cost
law for summation-type operations).  We also discuss the $\recip = 1\,\text{node}$
discovery as a dynamic application of these techniques.  The paper closes with
the open question: can $\addgen < 11\,\text{nodes}$?
\end{abstract}

%% =============================================================================
\section{Introduction: Four Complementary Lower-Bound Techniques}
\label{sec:intro}
%% =============================================================================

Every entry in the SuperBEST routing table has two components: an upper bound
(an explicit construction) and a lower bound (a proof that fewer nodes are
impossible).  Lower bounds are generally harder.  This note catalogues four
structural techniques that, together, cover all lower bounds in the SuperBEST
table and provide a framework for future entries.

\begin{definition}[SuperBEST routing table]
The SuperBEST table records, for each arithmetic operation $f$,
the value $\SB(f) = \min\{k : \exists\;\text{$k$-node $\F$-tree computing $f$}\}$.
Leaves are drawn from $\{x,y\}\cup\mathbb{Q}$ (rational-terminal convention).
\end{definition}

The four techniques are:
\begin{enumerate}
\item[\textbf{T35}] \textbf{Non-mergeability.}  Operator branches in an
  $\F$-tree cannot be ``merged'' unless their output types are algebraically
  compatible.  This rules out certain short trees by a structural incompatibility argument.
\item[\textbf{T36}] \textbf{Tree counting.}  Count the number of distinct
  $k$-node $\F$-trees and compare against the function's distinguishing points.
  If the function imposes more independent constraints than any $k$-node tree
  can satisfy, then $\SB(f) > k$.
\item[\textbf{T37}] \textbf{Disjointness.}  The value sets (images) of distinct
  operator branches are disjoint on relevant domains.  A function whose value
  set intersects two disjoint branches cannot be computed in a single node.
\item[\textbf{T40}] \textbf{Linear cost law.}  For summation-type operations
  $f(x_1, \ldots, x_m)$ that are linear in each argument, the cost satisfies
  $\SB(f) \ge C \cdot m$ for a constant $C > 0$.  This gives tight lower bounds
  for add and sub.
\end{enumerate}

The rest of this paper formalizes each technique and shows how they apply to
the SuperBEST table.

%% =============================================================================
\section{Techniques T35, T36, T37: Operator-Level Bounds}
\label{sec:operator}
%% =============================================================================

\subsection{T35: Non-Mergeability of Operator Branches}

\begin{definition}[Branch type]
For a node $v$ in an $\F$-tree computing $f$, the \emph{branch type} of $v$
is the ordered pair $(O, \sigma)$ where $O\in\F$ is the operator at $v$ and
$\sigma \in \{\mathrm{left}, \mathrm{right}\}$ indicates which child subtree
carries the free variable $x$ (or the dominant variable in a multi-variable tree).
\end{definition}

\begin{theorem}[T35: Non-Mergeability]
\label{thm:T35}
Let $f$ be a function with $\SB(f) = k$.  Suppose $f$ admits two distinct
decompositions
\[
  f(x) = O_1(A_1(x), B_1(x)) = O_2(A_2(x), B_2(x))
\]
where $O_1 \ne O_2$ are operators in $\F$ and $A_i, B_i$ are sub-expressions.
If the domains of $A_1$ and $A_2$ are disjoint (or the algebraic form of the
two decompositions is structurally incompatible), then no single $\F$-tree
of depth less than $k$ can realize both decompositions simultaneously.
In particular, if every 1-node $\F$-tree $O(a,b)$ fails to produce $f(x)$
for some witness point $(x_0,a_0,b_0)$, then $\SB(f) \ge 2$.
\end{theorem}

\begin{proof}
A 1-node tree is $O(\ell_1, \ell_2)$ for leaves $\ell_1, \ell_2$.  There
are finitely many choices of $(O, \ell_1, \ell_2)$ (16 operators times the
leaf set $\{x, y\} \cup \mathbb{Q}$).  For each such tree $T$, either
$T(x) = f(x)$ for all $x$ in the domain, or there exists a witness $x_0$
where $T(x_0) \ne f(x_0)$.  The non-mergeability argument provides a canonical
witness point for each operator family:
\begin{itemize}
  \item Class~I operators with $x$ first: take the derivative at any $x_0$;
    the derivative is exponential in $x$, not matching the required $f'(x_0)$.
  \item Class~II operators: evaluate at test points showing functional mismatch.
\end{itemize}
If all 16 operators fail at some witness, then $\SB(f) \ge 2$.  Inductively,
if all $(k-1)$-node trees fail, $\SB(f) \ge k$.
\end{proof}

\begin{remark}
T35 is the formalization of the ``derivative obstruction'' and ``slope barrier''
arguments used for $\negop$, $\mulop$, $\subop$, and $\divop$ in the audit paper.
The ``structural incompatibility'' formulation makes precise the intuition that
certain operator types cannot produce certain derivative patterns regardless
of what leaves are substituted.
\end{remark}

\subsection{T36: Tree Counting}

\begin{theorem}[T36: Tree-Counting Lower Bound]
\label{thm:T36}
Let $f$ be a function of $m$ real variables that is algebraically
independent of any $k$-node $\F$-tree.  Specifically, suppose there exist
$N > |\mathcal{T}_k|$ constraint points $(x_1, \ldots, x_N)$ such that
$f(x_i) = v_i$ for distinct target values $v_i$, and no $k$-node $\F$-tree
$T$ satisfies $T(x_i) = v_i$ for all $i$ simultaneously.
Then $\SB(f) \ge k+1$.
\end{theorem}

\begin{proof}
The set $\mathcal{T}_k$ of all $k$-node $\F$-trees is finite (there are at
most $16^k \cdot |\mathcal{L}|^{k+1}$ such trees for leaf set $\mathcal{L}$).
Each tree $T \in \mathcal{T}_k$ defines a real-analytic function.  If $f$
cannot be expressed as any such function (by identity on an open domain),
then $f \notin \{T : T \in \mathcal{T}_k\}$, hence $\SB(f) > k$.
\end{proof}

\begin{remark}
T36 provides a finite witness argument: it suffices to show that $f$ is not
in the (finite) list of $k$-node $\F$-functions.  For $k=0$, this reduces
to the ``leaf scan'' argument (no constant or identity leaf computes $f$).
For $k=1$, one checks all 16 operators with all leaf combinations.
In practice, the T35 derivative obstruction provides a single algebraic
certificate that covers all 16 operators simultaneously, making T36 redundant
in most cases.
\end{remark}

\subsection{T37: Disjointness of Value Sets}

\begin{theorem}[T37: Disjointness]
\label{thm:T37}
Let $S_O = \{O(\ell_1,\ell_2) : \ell_1,\ell_2 \in \mathcal{L}\}$ denote the
value set of operator $O \in \F$ applied to all leaf combinations.  If the
value sets $\{S_O\}_{O\in\F}$ are pairwise disjoint on a relevant domain $D$,
and $f|_D$ is not in any $S_O$, then $\SB(f) \ge 2$.
\end{theorem}

\begin{proof}
A 1-node tree is $O(\ell_1,\ell_2) \in S_O$ for some $O$.  If $f|_D \notin S_O$
for all $O$, then $f$ is not computed by any 1-node tree, giving $\SB(f) \ge 2$.
\end{proof}

\begin{remark}
T37 is most useful when the operators have non-overlapping structural forms,
such as comparing exponential-type operators (Class~I) with polynomial-type
target functions.  For instance, $\recip(x) = 1/x$ has a pole at $x=0$;
Class~I operators of the form $h(e^{\varepsilon x}, c)$ are bounded below by
a positive constant on $(0,\infty)$ (for appropriate $h$), so they cannot
equal $1/x$.  T37 formalizes this ``structural gap'' between operator types.
\end{remark}

\begin{corollary}[T37 for $\recip$]
Every Class~I operator applied to leaf inputs is bounded or unbounded in the
exponential-growth class; none has a simple pole at a positive real number.
The Class~II operator $\ELSb(0,x) = 1/x$ is the unique 1-node construction,
and it is valid because $\ELSb$ has the algebraic structure $e^a/b$ which
produces $1/x$ exactly when $a=0$.
\end{corollary}

%% =============================================================================
\section{T40: The Linear Cost Law for Summation-Type Operations}
\label{sec:T40}
%% =============================================================================

\begin{definition}[Summation-type operation]
A function $f:\mathbb{R}^m \to \mathbb{R}$ is \emph{summation-type} if it is
$\mathbb{R}$-linear in each argument separately:
\[
  f(\ldots, x_i + c, \ldots) = f(\ldots, x_i, \ldots) + c \cdot \partial_i f
  \quad\text{for all }c\in\mathbb{R},
\]
where $\partial_i f$ is the partial derivative with respect to $x_i$, which
must be nonzero and independent of $x_i$.
\end{definition}

\begin{theorem}[T40: Linear Cost Law]
\label{thm:T40}
Let $f:\mathbb{R}^m \to \mathbb{R}$ be summation-type with $\partial_i f \ne 0$
for all $i = 1, \ldots, m$.  Then
\[
  \SB(f) \ge \lceil m/2 \rceil.
\]
More precisely, each ``genuinely new'' linear dependence on an argument $x_i$
requires at least one $\F$-node to introduce, giving:
\[
  \SB(f) \ge \text{(number of distinct linear dependencies on distinct arguments)}/2.
\]
\end{theorem}

\begin{proof}[Proof sketch]
Each $\F$-node introduces at most one new ``linear slot'': Class~I nodes
introduce $e^{\varepsilon x}$ or $\varepsilon'\ln y$ (neither is linear in $x$ or $y$
for all values), while Class~II nodes $\ELAd(x,y) = e^x y$ and $\LEdiv(x,y) = x-\ln y$
introduce exactly one linear-type dependence ($y$ in $\ELAd$ when $e^x$ is treated
as a constant, and $x$ in $\LEdiv$).

For a summation-type $f$ with $m$ independent linear arguments, each argument
must be introduced by a separate path in the tree.  Since the tree is binary
and each node can route at most two subexpressions, at least $\lceil m/2\rceil$
nodes are needed to ``carry'' all $m$ arguments to the output with the correct
linear dependence.

Formally: any $\F$-tree $T$ computing $f$ has the property that the Jacobian
$(\partial T/\partial x_i)_{i=1}^m$ matches the constant vector
$(\partial_i f)_{i=1}^m$.  Each internal node of $T$ contributes at most two
nonzero partial derivative ``channels''.  For $m$ nonzero channels, we need
at least $\lceil m/2 \rceil$ nodes.
\end{proof}

\begin{example}[T40 applied to subtraction]
$\subop(x,y) = x - y$ is summation-type with $m=2$, $\partial_x = 1$,
$\partial_y = -1$.  By T40, $\SB(\subop) \ge 1$.  The structural lower bound
from T35 improves this to $\SB(\subop) \ge 2$ (no single node achieves both
$\partial_x=1$ and $\partial_y=-1$ simultaneously, as shown in the audit paper).
\end{example}

\begin{example}[T40 applied to general addition]
For $\addgen(x_1,\ldots,x_m) = x_1 + \cdots + x_m$ with $m$ arguments,
T40 gives $\SB(\addgen) \ge \lceil m/2\rceil$.  For $m=2$ (binary add),
this gives $\SB(\addop) \ge 1$, which T35/T37 strengthen to $\ge 3$.
For large $m$, T40 gives nontrivial bounds: $\SB(\addgen_m) \ge \lceil m/2\rceil$.
\end{example}

\begin{theorem}[T40 Corollary: Cost of $m$-ary sum]
\label{cor:T40msum}
For $m$-ary summation $\addgen_m(x_1,\ldots,x_m) = \sum_{i=1}^m x_i$,
\[
  \SB(\addgen_m) \ge \lceil m/2 \rceil.
\]
This bound is tight: a balanced binary tree of $\addop$ calls, each costing
$3$ nodes, gives $\SB(\addgen_m) \le 3\lceil \log_2 m \rceil$.
\end{theorem}

%% =============================================================================
\section{\texorpdfstring{$\recip = 1\,\text{node}$}{recip = 1 node}
  as Dynamic Extension}
\label{sec:recip}
%% =============================================================================

The discovery that $\recip(x) = 1/x$ requires only 1 node via $\ELSb(0,x)$
illustrates the ``dynamic'' aspect of the SuperBEST table: as new operators
are added to the family $\F$, previously optimal entries may be improved.

\begin{theorem}[Recip = 1 node: dynamic discovery]
\label{thm:recip-dynamic}
In the original 6-operator EML family $\mathcal{F}_6 = \{\EML, \ldots\}$
(without $\ELSb$), the minimum node count for $\recip$ was $\SB_6(\recip) = 2$.
After expanding to $\F = \mathcal{F}_{16}$ (adding $\ELSb$, $\ELAd$, $\LEAd$, $\LEdiv$),
we have $\SB_{16}(\recip) = 1$.
\end{theorem}

\begin{proof}
\emph{Lower bound in $\mathcal{F}_6$:}
The six original operators $\{\EML, \EML^-, \ldots\}$ all have the form
$e^{\varepsilon x} \pm \ln y$ or $e^{\varepsilon x} \cdot \ln y$.
For $O(x,c) = 1/x$ with constant $c$:
the derivative $O'(x)$ would need to equal $-1/x^2$ for all $x$.
All Class~I-form derivatives are either exponential in $x$ or of the form
$\pm 1/x$; none matches $-1/x^2$.
For $O(c,x) = e^{\varepsilon c}\ln^{\delta}(x)$: the only form matching
$1/x$ would be $e^0 / x = 1/x$, but no operator in $\mathcal{F}_6$ is of the
form $e^a/x$.  Hence $\SB_6(\recip) \ge 2$.

\emph{Upper bound in $\mathcal{F}_{16}$:}
$\ELSb(0,x) = e^0/x = 1/x$ (1 node).

The reduction from 2 to 1 node is the ``R16 dynamic discovery'': expanding
$\mathcal{F}_6$ to $\mathcal{F}_{16}$ by adding $\ELSb$ creates a new 1-node
path to $\recip$.
\end{proof}

\begin{remark}[Implications for T35/T37]
The recip discovery shows that T35 (non-mergeability) and T37 (disjointness)
are \emph{family-dependent}: a function that is not computable in 1 node within
$\mathcal{F}_6$ may become 1-node computable in a larger family.  The
lower-bound techniques are always relative to the given operator family.
This motivates the problem of finding the \emph{minimal family} for each target
function: the smallest $\mathcal{F} \subseteq \mathcal{F}_{16}$ such that
$\SB_{\mathcal{F}}(f) = 1$.
\end{remark}

%% =============================================================================
\section{Open Question: Can \texorpdfstring{$\addgen < 11\,\text{nodes}$}{add\_gen < 11 nodes}?}
\label{sec:open}
%% =============================================================================

The SuperBEST table records $\addgen(x,y) = x+y$ (for general signed real
inputs, not restricted to positive reals) as requiring 11 nodes.  The 3-node
construction works only when inputs are restricted to $(0,\infty)$, since it
relies on $\ln x$ and $\ln y$ being defined.

The best known lower bound for general real addition is the T08-A result:
$\SB(\addop) \ge 3$ on $(0,\infty)^2$, but for general $\addgen$ on $\mathbb{R}^2$,
the lower bound remains unsharpened to match the 11-node upper bound.

\begin{definition}[General addition]
$\addgen: \mathbb{R}^2 \to \mathbb{R}$, $(x,y) \mapsto x+y$, where $x,y$
range over all of $\mathbb{R}$ (including negative values).
\end{definition}

\begin{proposition}[Known bounds]
\label{prop:add-bounds}
\mbox{}
\begin{itemize}
  \item $\SB(\addop|_{(0,\infty)^2}) = 3$ (proved via $\LEAd(\EXL(0,x),y)$-type
    constructions and the T08-A lower bound).
  \item $\SB(\addgen) \le 11$ (upper bound via a piecewise construction that
    handles sign cases using $\F_{16}$ branching).
  \item $\SB(\addgen) \ge 3$ (from the positive-input lower bound, since
    $\addgen$ restricted to $(0,\infty)^2$ equals $\addop$).
\end{itemize}
\end{proposition}

\begin{remark}[The open gap: $3 \le \SB(\addgen) \le 11$]
The 8-node gap between 3 and 11 is the largest unresolved gap in the
SuperBEST table.  The difficulty is that handling negative inputs requires
sign detection, which $\F_{16}$ cannot perform in fewer than a certain number
of nodes (all operators in $\F_{16}$ are continuous and do not directly
implement a sign function).

\textbf{Open question:} Is $\SB(\addgen) < 11$?  More specifically:
\begin{enumerate}
  \item Can a 5--8 node construction handle all sign cases of $x+y$?
  \item Is there a structural lower bound greater than 3 for $\addgen$?
  \item Does T40 (linear cost law) provide a better lower bound when the
    domain includes sign changes?
\end{enumerate}
\end{remark}

The T37 disjointness technique suggests a possible approach: the value sets
of $\F_{16}$ operators have restricted domain structure (all involve
$e^{\pm x}$ or $\ln(\text{pos})$), and negative real outputs can only arise
from certain operator forms.  Characterizing which operators can produce
negative outputs might tighten the lower bound.

%% =============================================================================
\section{Conclusion}
\label{sec:conclusion}
%% =============================================================================

We have presented four complementary lower-bound techniques (T35--T37, T40)
that together form the structural foundation for the SuperBEST routing table.

\begin{table}[h]
\centering
\caption{Summary of lower-bound techniques and their applications.}
\label{tab:techniques}
\smallskip
\begin{tabular}{@{}llll@{}}
\toprule
\textbf{Technique} & \textbf{Mechanism} & \textbf{Key application} & \textbf{Strength} \\
\midrule
T35 (Non-mergeability) & Derivative class mismatch & $\negop$, $\mulop$, $\subop$, $\divop$ & Covers all 1-node failures \\
T36 (Tree counting) & Finite list exhaustion & $k$-node impossibility proofs & Brute-force completeness \\
T37 (Disjointness) & Value-set separation & $\recip$ in $\mathcal{F}_6$ vs.\ $\mathcal{F}_{16}$ & Domain-specific \\
T40 (Linear cost law) & Jacobian channel count & $\subop$, $\addgen_m$ & Multi-argument bounds \\
\midrule
\multicolumn{2}{l}{Dynamic extension: $\recip = 1\,\text{node}$}
  & R16-C1 ($\ELSb$ discovery)
  & Family-dependent \\
\midrule
\multicolumn{2}{l}{Open question: $\SB(\addgen) < 11\,\text{nodes}$?}
  & Gap: $3 \le \SB \le 11$
  & Unresolved \\
\bottomrule
\end{tabular}
\end{table}

The four techniques are not mutually exclusive; they are most effective when
applied in combination.  T35 rules out individual operators by derivative
mismatch.  T37 rules out entire operator classes by value-set separation.
T36 provides completeness at each depth level.  T40 gives depth-lower bounds
that scale with the number of arguments.  Together, they constitute a
systematic toolbox for proving structural optimality of $\F_{16}$-routing tables.

\bigskip
\noindent\textit{Almaguer, A.R.\ (2026).
Unified Structural Foundations of SuperBEST: Four Lower-Bound Techniques.
monogate research.\ \url{https://monogate.org}}

\end{document}
